% =====================================================================
% 決定力場力学 論文5 — ドラフト日本語版(v0.1, 2026-08-24)
% v0.1(第191便): 論文4/5 分割(原仮定者裁定 2026-08-24「論文4/5分割: 採用」)により新設。
%   旧論文4ドラフト(スケール横断の創発 v0.3)の**ミクロセクター**を引き取る:
%   - 「ミクロスケール: 内部スピン温度の熱化」章を**逐語移設**(本文・実測数値は不変 —
%     ハーネス exp-4-90/91/93)。
%   - 創発の標準(E0–E3 格付け+6標準試験)は本稿単独で読めるよう再掲(論文4と共有)。
%   - 隠れ質量ラダー(旧論文4マクロ節の骨子)は分割線
%     「論文4=重力・引きずり・形態(マクロ)/論文5=スピン・熱・光学(ミクロ)」に従い本稿へ。
%   新章(コアv2相転移・減光クラス D0–D6・減光/消失の設計制約実験)は**計画章の骨格**であり、
%   実測数値の主張はまだ無い。実在の恒星に関する主張ではない。
% 英語版: dfm-paper5.tex(正本 — 数値・構成は同一)
% =====================================================================
\documentclass[a4paper,10pt,twocolumn]{ltjsarticle}

\usepackage{luatexja-fontspec}
\usepackage[haranoaji]{luatexja-preset}
\usepackage{amsmath,amssymb}
\usepackage{graphicx}
\usepackage{booktabs}
\usepackage{hyperref}

\title{決定力場力学におけるミクロスケールの相変移:\\
内部スピン熱力学・コア転移・暗い天体の光学}
\author{(著者ブロックは論文1・2と同様 — TODO)}
\date{2026-08-24(ドラフト v0.1)}

\begin{document}
\maketitle

\section*{概要(骨格)}
% --- 主張すること ---
% 1. DFM トイ宇宙のミクロセクター(内部スピン温度 T_int・伝導チャネル・自転減光 lS_eff)を、
%    本シリーズ共通の段階格付け標準(E0–E3・6標準試験)で検証する。
% 2. ミクロ熱化(実測 — 旧論文4ドラフトから移設): 群間熱平衡化は伝導チャネルの真正な帰結
%    (全チャネル遮断で完全凍結)。減衰率 λ ≈ λ0 (κs/0.15)^{1.006} (n/120)^{1.084}。
%    平衡温度は C·m 重み台帳から予測可能(非等質量 2.09 — 単純平均 3.13 ではない)。
% 3. 計画: コアv2 相転移(収縮→lS_eff 飽和・パワーボール・空洞)/減光クラス D0–D6 の宣言的
%    定義/深い減光・消失様光度曲線の設計制約実験(実在の変光星・消えた星との同定はしない)。
% 4. 隠れ質量ラダー: 「暗さ」と「重さ」のモデル内分離。
% --- 主張しないこと ---
%    現実熱力学の説明・実在天体との同定・実在ダークマター・3D への定量外挿はしない。
概要 TODO。

\section{序論}
DFM は、全物質が源となる単一のスカラー場(決定度 $W$)が局所慣性系・重力・時計率・光学を
与える2次元の関係論的トイ宇宙である(論文1)。これは教育目的の計算的思考実験であって、
実宇宙の記述として提案するものではなく、以下のどの現象も実在の天体に関する主張ではない
(限界節)。論文1が基礎、論文2が箱宇宙、論文3が「較正一致≠機構同定」の哲学による現実較正、
論文4がマクロ側(重力・引きずり・形態)の天体再現を担う。本稿は同じ計画の\emph{ミクロ側} —
内部スピン温度の熱力学・コアモデルの相転移・暗い天体の光学(減光ラダー)— を収める。
論文4/5の分割線(2026-08-24 採用)は機構に基づく: 論文4=重力・引きずり・形態(マクロ)/
論文5=スピン・熱・光学(ミクロ)。実測済みの章(ミクロ熱化)は旧合同ドラフトから逐語移設し、
残りの章は測る前にゲートを宣言した計画である。

\section{創発の標準(方法論の中核 — 論文4と共有)}
% 表1: E0–E3 の定義。6標準試験の定義と運用規則(秩序変数の事前宣言・
% 負の結果も同一書式で記録)。旧合同ドラフトで導入 — 論文4と逐語共有。
TODO。

\section{ミクロスケール: 内部スピン温度の熱化}
% 実測(exp-4-90/91/93): 全断ノックアウトで λ=0(R=1.0000 の凍結)/
% κs 単独遮断でも λ≈0(群間輸送はほぼ伝導のみ)/ κs 3値で λ ほぼ比例/
% 16seed で λ=2.45e-5〜4.20e-5(ばらつき1.71倍)/ n 3段階で単調増/
% スケーリング fit: λ ≈ λ0 (κs/0.15)^{1.006} (n/120)^{1.084}, λ0=3.28e-5/步/
% 平衡時間 ≈6.8万步(時間窓の成果 — 初版 t_eq 判定の失敗と λ への差し替えの経緯を明記)/
% 摂動回復: 温度再分裂 λ2/λ1=0.96〜1.20・速度キック λ3/λ1=1.10〜1.32/
% 平衡温度予測: T_eq=Σ C·m·T₀/Σ C·m — 等質量 3.13・非等質量 2.09(質量重み)を確証。
% 副産物: kRep>0 では熱圧の仕事が T_int を差し引く(E5′×tint の設計 — 熱⇄仕事チャネル)。
% [第100便: 本文起こし — 英語版 dfm-paper4.tex(旧 dfm-paper3.tex v0.2)の対応節と数値・構成同一。
%  G=0 の分子系のため c₀=30 再パラメータ化(第96〜97便)の影響なし(力学 bit 不変)]

\subsection{系と秩序変数}

ミクロスケールの試験系は意図的に最小である: 閉じた箱に同一の粒
$120+120$ 個を、左を低温($T_{\mathrm{int}}=0.01$)・右を高温
($T_{\mathrm{int}}=6.25$)として用意し、重力と引きずりは切る
($G=0$, $k_{\mathrm{frame}}=0$)。温度は A9$''$ 改訂の内部状態変数
$T_{\mathrm{int}}$(熱エネルギー $E=C\,m\,T_{\mathrm{int}}$)であり、
ここでの「熱平衡化」とは伝導チャネル E10$'$ と衝突による
$T_{\mathrm{int}}$ の輸送を指す — 並進運動の等分配ではない。
宣言する秩序変数は群間コントラスト $R(t)=\Delta T(t)/\Delta T(0)$、
判定量はその減衰率 $\lambda$($\ln R(t)$ の最小二乗傾き)である。

この判定量の選択自体が、記録された差し替えである。ハーネスの初版は
平衡化時間 $t_{\mathrm{eq}}$ を 3000 步の観測窓で判定し、
\emph{全条件で失敗した}: この系の伝導時間尺度は長く、実測減衰の外挿で
$R=0.1$ 到達は約 $6.8\times10^{4}$ 步にあたる。窓を試験が通るまで
広げるのではなく、判定量を率 $\lambda$ に置き換え、失敗した初版の
経緯をハーネス冒頭(\texttt{tests/exp-4-90.mjs})に記録し、
時間尺度そのものは標準試験の時間窓試験(試験6)の結果として残した。

\subsection{6つの標準試験}

\emph{ノックアウト(試験1)。}熱チャネルを全て切る
($\kappa_s=\mu_F=\gamma_n=0$)とコントラストは厳密に凍結する:
3 seed で $R=1.0000$・$\lambda=0$。伝導\emph{だけ}を切る
($\kappa_s=0$・衝突は残す)と率はノイズ水準
($|\lambda|\lesssim 2\times10^{-6}$/步・符号すら不安定)へ潰れる:
群間輸送は実質すべて伝導チャネルが担い、衝突は各群の内部で散逸する
だけである。現象を支える機構はちょうど1つで、それを抜けば時計が止まる。

\emph{用量反応(試験2)。}$\kappa_s \in \{0.05, 0.15, 0.45\}$ の掃引で
$\lambda$ は単調・ほぼ比例で応答する(下の大域フィットの指数は 1.006)。

\emph{多seed(試験3)。}16 seed で $\lambda = 2.45\times10^{-5}$〜
$4.20\times10^{-5}$/步(ばらつき 1.71 倍・中央値 $3.28\times10^{-5}$)。
緩和しない seed も、チャネルなしで緩和する seed もない。

\emph{粒子数スケーリング(試験4)。}片側 $n=60/120/240$ で率は単調に
増える。用量掃引と合わせて
\begin{equation}
\lambda \;\approx\; \lambda_0
\left(\frac{\kappa_s}{0.15}\right)^{1.006}
\left(\frac{n}{120}\right)^{1.084},
\qquad \lambda_0 = 3.28\times10^{-5}\ \text{/步}
\label{eq:microfit-ja}
\end{equation}
に載る — 伝導係数と粒子密度への比例が数\%で成り立つ
(\texttt{tests/exp-4-91.mjs})。

\emph{摂動回復(試験5)。}平衡化の進行後に2通りの擾乱を与える:
(a)温度の再分裂は再緩和率比 $\lambda_2/\lambda_1 = 0.96/1.15/1.20$
(seed 別)、(b)アプリ正式の運動量保存injectorによる速度キックは
$\lambda_3/\lambda_1 = 1.10/1.32/1.32$。いずれも宣言済みの受理帯
$[0.5, 2]$ に収まる: この緩和は初期条件の一回性ではなく、
機構が中断のたびに立て直す率である。

\emph{時間窓(試験6)。}正直な数字は初版ハーネスを壊した当のもので
ある: 既定結合での完全平衡化は約 $7\times10^{4}$ 步かかる。
率を報告し、時間尺度が長いことも報告する。

\subsection{台帳水準の予測}

上の試験群は頑健性を立てるが、セクターがモデル内で\emph{定量的}で
あることまでは示さない。そこで伝導のみ構成
($\mu_F=\gamma_n=k_{\mathrm{rep}}=0$)を使う: 熱
$Q=\sum C\,m\,T_{\mathrm{int}}$ は厳密に保存され、平衡温度は
$C\,m$ 重み平均
\begin{equation}
T_{\mathrm{eq}} \;=\; \frac{\sum_i C_i m_i T_{0,i}}{\sum_i C_i m_i}
\label{eq:teq-ja}
\end{equation}
になるはずである。実測(\texttt{tests/exp-4-93.mjs}・各2 seed):
等質量では予測 $T_{\mathrm{eq}}=3.13$ に両群が上下から単調に接近し、
$Q$ のドリフトは相対 $10^{-6}$ 台。判別力があるのは非等質量
(左群 $m=2$)で、素朴な算術平均なら $3.13$ のままだが台帳は
$T_{\mathrm{eq}}=2.09$ を予測する — 実測の両群は $2.09$ を挟んで
接近し、$Q$ は同じく $10^{-6}$ 台で保存される。この予測は非自明で
あり(質量重みと素朴平均を区別する)、アプリの保存量モニタが表示する
のと同じ簿記から導かれる。既定の散逸あり構成では同じ台帳に衝突散逸の
項が加わる — 明確な主張は、検証対象のチャネルだけが開いた構成で
意図的に行う。

6試験と台帳予測により、ミクロセクターは第2節の標準試験一式を
満たす。立つ主張は控えめで精密である: このトイ宇宙において、群間の
熱平衡化は宣言された1チャネルの真正で頑健な帰結であり、その率則
[式(\ref{eq:microfit-ja})]と平衡点[式(\ref{eq:teq-ja})]は
seed・密度・摂動を生き延びる。

\section{コアv2 相転移(計画)}
% 計画章 — 実測数値の主張はまだ無い。収縮下のコアv2状態変数(M_c, R_c, J_core)・
% 実効減光 lS_eff = min(1, |s|R/c_surf) の飽和を転移マーカーとする・パワーボール/空洞領域・
% 測る前に宣言する秩序変数。ゲートは6標準試験+保存則帳簿+dt/軟化長収束。
TODO。

\section{減光ラダー: 宣言的光学クラス D0--D6(計画)}
% 計画章。D0–D6 はトイ宇宙の減光機構の**宣言されたクラス**(「明るい」から
% 「力学的に重く光学的に暗い」まで)であり、実在天体の分類ではない。成果物はクラス定義・
% 機械検査可能なシグネチャ・否定対照(減光は力学とビット厳密に分離 — 論文4メゾ章で確立済み)。
TODO。

\section{設計制約実験: 深い減光と消失(計画)}
% 計画章。宣言済み光学クラスのどのパラメータ領域が (i) 深く不規則な減光光度曲線
% (ii) 消失様光度曲線をトイ宇宙内で作り得るかを問う設計制約実験。これはモデル自身の
% クラスへの設計制約であり、実在の変光星・消えた星との同定では**ない**。
% 論文4のマクロ本線からは独立(相互にゲートしない)。
TODO。

\section{隠れ質量ラダー}
% 旧論文4マクロ節の骨子を分割線(光学・簿記=ミクロ側)に従い移設:
% 真の質量/光度プロキシ/ニュートン力学質量 v²r/G/引きずり超過・kFrame=0 対照。
% 「暗さ」と「重さ」のモデル内分離。
TODO。

\section{限界と負の主張}
% 固定リスト(論文1〜4の精神を継承):
%  1. 2D のみ・3D への定量外挿なし。
%  2. Lorentz 不変なし・絶対同時性・因果伝播なし。
%  3. 減光はモデル固有(輻射輸送ではない)。
%  4. 熱セクターはアナロジー(簿記はモデル内部のもの)。
%  5. 減光クラス D0–D6 はトイ宇宙の宣言クラス — 実在の変光星・消えた星との同定はしない。
%  6. ダークローター≠ダークマター同定(論文4参照)。
%  7. 存在量・形成率の天体物理的主張はしない。
TODO。

\section{再現性}
% 現象→プリセット→ハーネス→seed集合→時間窓→コミット→結果JSON の台帳表。
% 実測章のハーネス: exp-4-90(ミクロ標準試験)・exp-4-91(摂動回復+スケーリングfit)・
% exp-4-93(平衡温度予測)。計画章は測る前に同書式でハーネスを宣言する。
TODO。

\begin{thebibliography}{9}
\bibitem{paper1} 論文1(TODO)
\bibitem{paper2} 論文2(TODO)
\bibitem{paper3} 論文3(TODO)
\bibitem{paper4} 論文4(天体再現 — TODO)
\end{thebibliography}

\end{document}
